% ---------------------------------------------------------------------------
% Introduction draft -- filament CMB lensing with BOSS CMASS pairs
%
% Assumes natbib (\citet / \citep) and the macro
%   \newcommand{\hinvMpc}{h^{-1}\,\mathrm{Mpc}}
%
% Placeholder bib keys used below (rename to match your .bib):
%   Zeldovich1970, BondKofmanPogosyan1996, Cautun2014, FukugitaHoganPeebles1998,
%   Shull2012, CenOstriker1999, deGraaff2019, Tanimura2020, Nicastro2018,
%   Bartelmann2001, LewisChallinor2006, Colberg2005, Clampitt2016,
%   EppsHudson2017, Kondo2020, White2011, Jenkins1998, Planck2018lensing,
%   He2018, Chen2016, Sousbie2011, Tanimura2019, Klypin2016, Reid2016
%
% Numbers marked "% CHECK" should be verified against the source before submission.
% ---------------------------------------------------------------------------

\section{Introduction}
\label{sec:intro}

On scales of tens of megaparsecs the matter distribution of the Universe is
organised into a cosmic web of nodes, filaments, sheets, and voids. This
morphology is a generic outcome of anisotropic gravitational collapse: in the
Zel'dovich picture a density perturbation contracts first along one axis to
form a sheet, then along a second to form a filament, and only finally along
the third to form a compact node \citep{Zeldovich1970, BondKofmanPogosyan1996}.
Filaments are the dominant reservoir of matter in this network. In $N$-body
simulations they contain roughly half of the mass of the Universe while
occupying only a few per cent of its volume \citep{Cautun2014}. They are also
the natural home of the baryons that are not accounted for in stars, cold
gas, and the hot atmospheres of groups and clusters: at low redshift roughly
30 per cent of the cosmic baryon budget is unobserved
\citep{FukugitaHoganPeebles1998, Shull2012}, and hydrodynamic simulations
predict that most of it resides in a warm--hot intergalactic medium at
$T\sim10^{5}$--$10^{7}\,\mathrm{K}$ tracing the filamentary web
\citep{CenOstriker1999}. Measuring the matter content of filaments therefore
tests both the gravitational growth of structure and the census of baryons
\citep{deGraaff2019, Tanimura2020, Nicastro2018}.

Gravitational lensing is the most direct probe of this matter. The lensing
convergence is a weighted projection of the matter overdensity $\delta$ along
the line of sight,
\begin{equation}
  \kappa(\hat{\bm n}) = \int_{0}^{\chi_{\rm s}} \mathrm{d}\chi\,
  W(\chi)\,\delta(\chi\hat{\bm n},\chi), \qquad
  W(\chi) = \frac{3\Omega_{\rm m}H_{0}^{2}}{2c^{2}}\,
  \frac{\chi}{a(\chi)}\,\frac{\chi_{\rm s}-\chi}{\chi_{\rm s}},
  \label{eq:kappa}
\end{equation}
where $\chi$ is comoving distance and $\chi_{\rm s}$ the distance to the
source plane \citep{Bartelmann2001}. Because $\kappa$ responds to the total
matter density, a lensing measurement of a filament requires no assumption
about how galaxies or gas trace the dark matter. Equation~(\ref{eq:kappa}) also
sets the scale of the expected signal. For a source at the last-scattering
surface, $z_{\rm s}\simeq1090$, the kernel $W(\chi)$ is broad and peaks near
$z\sim2$ \citep{LewisChallinor2006}; at the redshift of the BOSS CMASS sample,
$z\simeq0.55$, its value is $W\simeq7\times10^{-5}\,\mathrm{Mpc}^{-1}$.
A filament of core overdensity $\delta\sim3$--$5$ and diameter of a few
megaparsecs, seen in projection, then produces
$\kappa\sim(0.5$--$1.5)\times10^{-3}$. A background galaxy population at
$z_{\rm s}\sim1$ has a kernel that peaks close to the lens redshift and
therefore yields a larger per-object signal, at the cost of shape noise and a
much smaller footprint.

Filaments cannot be located individually in a noise-dominated lensing map, so
all measurements to date stack many candidate structures. The candidates are
either identified algorithmically from the galaxy distribution or, as in this
work, defined by pairs of massive galaxies whose connecting axis fixes the
expected orientation of a bridge. Simulations show that pairs of
group- and cluster-mass haloes are connected by a filament with a probability
that decreases with separation, from near unity below $\sim5\,\hinvMpc$ to
roughly one third at $15$--$20\,\hinvMpc$ \citep{Colberg2005}. % CHECK fractions against Colberg+2005
CMASS galaxies occupy haloes of $M\sim10^{13}\,h^{-1}\mathrm{M}_{\odot}$
\citep{White2011}, so a pair at these separations is a plausible filament
anchor. The mean convergence around such a pair decomposes into three
pieces: the one-halo profiles of the two members, the isotropic two-halo
contribution from the correlated large-scale environment in which the pair
sits, and an anisotropic excess concentrated along the axis joining them. In
the language of correlation functions the isotropic terms are set by the
galaxy--matter two-point function, while the anisotropic excess is the
connected part of the galaxy--galaxy--matter three-point function
\citep{Clampitt2016}. It is this last term that we identify as the filament
signal, and it is a quantity with a definite $\Lambda$CDM prediction. Because
the connection probability falls with separation and the bridge density
declines as it lengthens, the amplitude of the excess should decrease with
pair separation. Measuring that dependence is a primary goal of this work.

Two selection effects govern the purity of a pair sample. The first is
projection: pairs that are close on the sky but distant along the line of
sight are not connected, and their inclusion dilutes the mean signal. The
second is redshift-space distortion. The observed line-of-sight separation
of two galaxies includes their relative peculiar velocity, and for haloes at
these separations the pairwise velocity dispersion is $\sigma_{v}\simeq500\,
\mathrm{km\,s^{-1}}$ \citep{Jenkins1998, deGraaff2019}. At $z=0.55$ this
corresponds to a comoving displacement $(1+z)\sigma_{v}/H(z)\simeq
6\,\hinvMpc$, which is why previous pair-stacking analyses have adopted a
line-of-sight cut of $5$--$6\,\hinvMpc$. Such a cut rejects some genuinely
connected pairs whose members have large relative velocities and admits some
unconnected pairs that happen to lie within the velocity window, but it
retains pairs whose axes are close to the plane of the sky, which is the
geometry in which the bridge is seen with least projection.

\subsection{Previous stacked lensing measurements of filaments}

Existing measurements divide according to whether the lensing source is a
population of background galaxies or the CMB.

\citet{Clampitt2016}, \citet{EppsHudson2017}, and \citet{Kondo2020} use galaxy
shapes. \citet{Clampitt2016} stacked the shear field around pairs of SDSS
luminous red galaxies at $z\approx0.25$ with projected separations of
$6$--$14\,\hinvMpc$ and line-of-sight separations below $6\,\hinvMpc$, using
SDSS source shapes over approximately $8000\,\mathrm{deg}^{2}$, and reported
a $4.5\sigma$ detection. Their four-point estimator combines the shear at
positions rotated by $90^{\circ}$ about each halo centre, which cancels the
tangential shear of two spherically symmetric haloes analytically and
isolates the bridge. \citet{EppsHudson2017} stacked CFHTLenS convergence
around BOSS LOWZ and CMASS pairs at a mean redshift $\langle z\rangle\approx
0.42$. They modelled the two-halo contribution with the stack of
non-physical pairs, selected with the same projected separation but large
line-of-sight separation, and reported a $5\sigma$ excess after subtracting
it from the physical-pair stack. \citet{Kondo2020} paired BOSS CMASS galaxies
at $z\approx0.55$ with first-year Subaru Hyper Suprime-Cam shapes. Although
the overlap footprint is only $\sim140\,\mathrm{deg}^{2}$ and contains
$14{,}422$ CMASS galaxies, the source density of $\sim22\,\mathrm{arcmin}^{-2}$
yielded a $3.9\sigma$ detection. Their mock analysis showed that intrinsic
variation among filaments together with projected large-scale structure
contributes approximately 60 per cent of the covariance, with source shape
noise supplying the remainder, so that deeper imaging alone would not
substantially improve the measurement.

The CMB provides a complementary source. Its lensing kernel is broad, the
reconstruction is free of galaxy shape noise and intrinsic alignments, and
the Planck maps cover most of the sky, so the full BOSS footprint can be
used. Against this, the Planck reconstruction has per-mode signal-to-noise
ratio of order unity only at multipoles $L\approx40$--$60$ and is strongly
noise dominated on the arcminute scales relevant to filaments
\citep{Planck2018lensing}, and it carries possible foreground-related
systematics. Two studies have stacked or cross-correlated Planck
convergence with filaments identified directly from the spectroscopic galaxy
distribution. \citet{He2018} applied a density-ridge algorithm to the CMASS
distribution \citep{Chen2016}, constructed a filament intensity map from the
resulting catalogue, and cross-correlated it with the Planck convergence map,
obtaining a $5\sigma$ detection and a filament bias $b_{\rm f}\approx1.5$.
\citet{Tanimura2020} identified $24{,}544$ filaments of length
$30$--$100\,\mathrm{Mpc}$ in LOWZ+CMASS with DisPerSE \citep{Sousbie2011},
masked the critical points at their ends, and stacked Planck $\kappa$ on the
filament spines, reaching $8.1\sigma$. These significances are not directly
comparable to pair-stacking results: the structures are longer, are selected
to be genuine overdensities, and are analysed with different estimators,
whereas a pair sample necessarily includes unconnected pairs that dilute the
mean.

The study of \citet{deGraaff2019} is the closest predecessor of this work.
Their primary measurement was the thermal Sunyaev--Zel'dovich signal of
filament gas, and they applied the same pipeline to the Planck convergence
map to estimate the filament matter density. They stacked around
$1{,}002{,}334$ pairs % CHECK: draft had 1,020,334; lit notes have 1,002,334
drawn from the BOSS CMASS catalogue used here, with transverse comoving
separations of $6$--$14\,\hinvMpc$ and line-of-sight separations below
$5\,\hinvMpc$. Each pair is rotated to a common orientation, rescaled to unit
separation, projected onto a rectangular grid, and reflected to exchange the
two members. Halo contributions are removed by fitting an isotropic radial
profile within a $60^{\circ}$ cone perpendicular to the pair axis and
subtracting the resulting two-dimensional model. Because the convergence map
was matched to the $10$-arcmin resolution of the Compton-$y$ map, the
lensing analysis was smoothed more heavily than the reconstruction requires.
They measured a mean filament convergence of
$\bar{\kappa}=(0.58\pm0.31)\times10^{-3}$ between the pair members, a
$1.9\sigma$ result, rising to $3.1\sigma$ when the region beyond the pair
is included.

\subsection{This work}

We revisit the combination of BOSS CMASS pairs with Planck CMB lensing,
following \citet{deGraaff2019}, with four changes in design.

First, we resolve the measurement in pair separation rather than integrating
over it. Where \citet{deGraaff2019} rescale every pair to a normalised
separation before stacking, we stack on a fixed comoving grid of
$101\times101$ pixels spanning $100\times100\,(\hinvMpc)^{2}$, using disjoint
separation bins centred at $5$, $10$, and $20\,\hinvMpc$. The finite width of
each bin smears the bridge modestly, but the construction allows us to test
the predicted decline of the anisotropic excess with separation described
above.

Second, we model the entire measurement in simulation. We construct a
CMASS-like mock catalogue from the BigMDPL $N$-body simulation
\citep{Klypin2016} and apply matched pair selection, map smoothing, stacking,
and halo subtraction to convergence maps generated from the same volume. This
supplies a separation-dependent $\Lambda$CDM reference for both the amplitude
and the transverse profile of the bridge, which can be compared directly with
the observed stacks.

Third, we use an independent halo-subtraction method. Where
\citet{deGraaff2019} fit an isotropic profile to the observed pair stack
itself, we subtract a control map built by displacing the stack of single
CMASS galaxies to $\pm d/2$ along the pair axis, in the spirit of the control
stacks of \citet{EppsHudson2017}. The two approaches fail in different ways.
Profile fitting can absorb part of the bridge into the halo model. The
control map, by construction, contains the one-halo term and the isotropic
two-halo term of an average CMASS galaxy, but galaxies selected to have a
close companion reside in denser-than-average environments, so the control
underestimates the isotropic two-halo contribution of paired galaxies and
leaves a positive residual that is not filamentary. We quantify this bias
with the mock catalogue, where the same control construction can be applied
to a simulation with known halo and filament content. % CHECK: confirm this test is in the analysis

Fourth, we vary the line-of-sight cut rather than adopting the
$5$--$6\,\hinvMpc$ value used in previous pair-stacking studies. Larger cuts
increase the pair count but admit more unconnected pairs, and, as discussed
above, peculiar velocities of $\sim500\,\mathrm{km\,s^{-1}}$ move pairs across
any threshold of a few $\hinvMpc$. We therefore measure directly how the cut
trades contamination against sample size, again using the mock to calibrate
the fraction of connected pairs at each cut.

We use the Planck 2018 minimum-variance lensing reconstruction
\citep{Planck2018lensing} synthesised at $N_{\rm side}=2048$ and smoothed with
an $8$-arcmin Gaussian kernel, compared with the $10$-arcmin smoothing of
\citet{deGraaff2019}. At $z=0.55$ the comoving distance is
$\chi\simeq1400\,\hinvMpc$, so $8$ arcmin corresponds to $\simeq3.3\,\hinvMpc$,
comparable to the expected filament core radius, while pair separations of
$5$, $10$, and $20\,\hinvMpc$ subtend approximately $12$, $24$, and $48$ arcmin.
The smallest separation bin is therefore only marginally resolved after
smoothing, and we treat it as a consistency check on the larger bins rather
than as an independent profile measurement. Because the reconstruction noise
exceeds the signal on all of these scales, the difference between $8$- and
$10$-arcmin smoothing has a modest effect on the recovered significance.

Throughout we adopt the Planck 2018 cosmology with $h=0.6766$ and quote
comoving distances in $\hinvMpc$. A positive filament signal denotes an excess
of convergence, and hence of projected matter, between the pair members
relative to the control.

The paper is organised as follows. Section~\ref{sec:data} describes the
CMASS and random catalogues, the Planck lensing map, and the BigMDPL mock
catalogue with its simulated convergence map. Section~\ref{sec:method}
presents the pair selection, stacking, control construction, band statistic,
and error estimation, applied identically to data and mock.
Section~\ref{sec:sim_expect} gives the $\Lambda$CDM expectation from the
mock. Section~\ref{sec:results} reports the observed stacks, the bridge
convergence at each separation, the comparison with the mock prediction and
with previous work, and the robustness tests. Section~\ref{sec:discussion}
summarises the results and discusses their implications.
